\chapter{Sistem Persamaan Linear, Determinan, dan Invers}

Sebuah penilai ancaman menjumlahkan tiga fitur dengan bobot yang belum diketahui,
$s=w_{1}f_{1}+w_{2}f_{2}+w_{3}f_{3}$. Kalibrasi pada tiga contoh berlabel memberi tiga
persamaan linear sekaligus:
\[
  2w_{1}+w_{2}-w_{3}=8,\qquad
  -3w_{1}-w_{2}+2w_{3}=-11,\qquad
  -2w_{1}+w_{2}+2w_{3}=-3 .
\]
Jawabannya $w=(2,3,-1)$, dan tidak ada yang mengejutkan.

Kemudian matriks koefisiennya diganti oleh hasil pengukuran lain,
\[
  S=\begin{pmatrix}1&2&3\\ 4&5&6\\ 7&8&9{,}001\end{pmatrix}.
\]
Baris kedua dikurangi baris pertama memberi $(3,3,3)$, sedangkan baris ketiga dikurangi
baris kedua memberi $(3,3,3{,}001)$: ketiga baris hampir bergantung. Bobot $w=(1,1,1)$
menghasilkan ruas kanan $(6,\,15,\,24{,}001)$, jadi sistem itu konsisten dan jawabannya tepat
$w=(1,1,1)$. Naikkan pengukuran ketiga sebesar $0{,}001$, dan jawabannya berpindah ke
$(2,-1,2)$. Perubahan relatif pada data $3{,}46\times10^{-5}$ menghasilkan perubahan relatif
pada jawaban $1{,}41$.

Ada empat pertanyaan, bukan satu: apakah sistem punya jawaban, berapa banyak, angka apa yang
memutuskannya tanpa menyelesaikan apa pun, dan seberapa jauh jawaban bergeser ketika data
bergeser. Bab 7 sudah memberi rank dan nullity untuk yang pertama. Yang belum ada tiga:
prosedur yang tercatat dan dapat diperiksa orang lain, satu bilangan yang memutuskan
invertibilitas sekaligus memberi rumus jawaban, dan ukuran kepekaan; ketiganya isi bab ini.

Determinan lazim diperkenalkan sebagai rumus kofaktor, lalu sifatnya dibuktikan satu per
satu. Kita balik: definisi dulu lewat tiga sifat, ketunggalan dibuktikan, rumus kofaktor
muncul sebagai akibat. Alasannya bukan gaya: rumus kofaktor menyembunyikan bahwa determinan
adalah satu-satunya fungsi bergantian yang ternormalisasi, dan fakta itulah yang membuat
bukti di bab ini pendek.

\begin{untukanak}
Tiga kawan menimbang tiga kotak dengan tiga timbangan berbeda. Setiap kawan menuliskan tiga
angka, dan dari angka-angka itu kita ingin tahu berat setiap kotak. Selama ketiga timbangan
mengukur hal yang berbeda, jawabannya mudah: setiap angka baru memberi satu keterangan baru.

Masalahnya begini. Timbangan kedua ternyata hampir menyalin timbangan pertama. Angkanya
tidak sama persis, hanya berbeda sedikit, mungkin karena debu di piringan. Timbangan ketiga
juga hampir menyalin keduanya. Sekarang ketiga keterangan itu hampir sama isinya. Kita
merasa punya tiga keterangan, padahal yang benar-benar baru hanya satu atau dua.

Di sinilah keanehan mulai. Kalau salah satu angka berubah seperseribu, berat kotak yang
kita hitung bisa berubah hampir setengah. Bukan karena hitungannya salah, tetapi karena
keterangannya memang hampir kembar. Ada satu bilangan yang dapat memberi tahu lebih dulu
apakah hal itu akan terjadi, sebelum kita menghitung apa pun.

Pernahkah kamu menanyakan hal yang sama kepada tiga orang, lalu mendapat tiga jawaban yang
isinya mirip, dan tetap tidak yakin? Di bab ini kita akan tahu mengapa kemiripan jawaban
justru melemahkan keterangan, dan bagaimana mengenalinya dari angkanya.
\end{untukanak}

\section{Eliminasi Gauss dan tiga nasib sebuah sistem}
\tingkat{dasar}

Sistem linear ditulis sebagai satu kesatuan, bukan sebagai daftar persamaan yang terpisah.
Caranya dengan menempelkan kolom ruas kanan ke matriks koefisien. Penyelesaian dikerjakan
pada susunan itu, dan setiap langkahnya adalah penggantian satu sistem oleh sistem lain yang
himpunan jawabannya sama.

\kamus{eliminasi}{elimination}{Mengganti pertanyaan dengan pertanyaan yang lebih sederhana
tanpa mengubah jawabannya, seperti menyederhanakan pecahan tanpa mengubah nilainya.}

\begin{definisi}[Matriks diperbesar dan operasi baris elementer]
Untuk sistem $A\mathbf{x}=\mathbf{b}$ dengan $A$ berukuran $m\times n$, \emph{matriks
diperbesar} adalah matriks $\bigl(A\mid\mathbf{b}\bigr)$ berukuran $m\times(n+1)$. Tiga
\emph{operasi baris elementer} adalah: menukar dua baris; mengalikan satu baris dengan
skalar tak nol; menambahkan kelipatan satu baris pada baris lain. Dua matriks disebut
\emph{ekuivalen baris} bila yang satu dapat dicapai dari yang lain dengan operasi-operasi
itu.\index{eliminasi}
\end{definisi}

Operasi itu dipilih karena masing-masing dapat dibatalkan. Menukar dua baris membatalkan
dirinya sendiri, mengalikan dengan $c\neq0$ dibatalkan dengan mengalikan $1/c$, dan
menambahkan $c$ kali baris $i$ pada baris $j$ dibatalkan dengan mengurangi $c$ kali baris
$i$. Karena itu tidak ada informasi yang hilang.

\begin{teorema}
Bila dua sistem linear mempunyai matriks diperbesar yang ekuivalen baris, kedua sistem itu
mempunyai himpunan jawaban yang sama.
\end{teorema}

\begin{proof}
Cukup diperiksa untuk satu operasi, karena operasi berulang mengikuti. Misalkan $S$ adalah
sistem semula dan $S'$ sistem setelah satu operasi. Setiap persamaan $S'$ adalah kombinasi
linear persamaan-persamaan $S$, sehingga setiap jawaban $S$ juga menjawab $S'$. Sebaliknya,
karena operasi itu dapat dibatalkan oleh operasi sejenis, setiap persamaan $S$ adalah
kombinasi linear persamaan $S'$, sehingga setiap jawaban $S'$ juga menjawab $S$. Kedua
himpunan jawaban karena itu saling memuat.
\end{proof}

Bentuk yang dituju adalah bentuk tangga. Setiap baris tak nol dimulai lebih ke kanan
daripada baris tak nol di atasnya, dan semua entri di bawah posisi pertama itu nol. Kolom
tempat angka pertama sebuah baris tak nol berada disebut kolom pivot, dan angka itu sendiri
disebut pivot.

\begin{definisi}[Bentuk eselon baris]
Matriks disebut dalam \emph{bentuk eselon baris} bila semua baris nol berada di bawah, dan
entri tak nol pertama setiap baris tak nol terletak lebih kanan daripada entri tak nol
pertama baris di atasnya. Bentuk itu disebut \emph{tereduksi} bila setiap pivot bernilai
$1$ dan merupakan satu-satunya entri tak nol pada kolomnya.\index{eselon}\index{pivot}
\end{definisi}

Bentuk tereduksi itu tunggal. Fakta ini dipakai setiap kali kita berbicara tentang ``rank
yang dihitung dengan eliminasi'', dan buktinya tidak sependek kelihatannya.

\begin{teorema}[Ketunggalan bentuk eselon tereduksi]
Setiap matriks ekuivalen baris dengan tepat satu matriks dalam bentuk eselon baris
tereduksi.
\end{teorema}

\begin{proof}
Bukti dilewati di sini. Pernyataan ini dibuktikan lengkap di Hoffman dan Kunze,
\emph{Linear Algebra} edisi ke-2, bagian 1.4, dengan menunjukkan bahwa ruang jawaban
$A\mathbf{x}=\mathbf{0}$ menentukan letak pivot secara tunggal. Yang kita pakai di bab ini
hanya akibatnya: banyaknya pivot dan letaknya tidak bergantung pada urutan operasi yang
dipilih.
\end{proof}

Dengan bentuk tereduksi, seluruh pertanyaan tentang jumlah jawaban dapat dijawab dengan
menghitung pivot. Inilah tiga nasib sebuah sistem.

\begin{teorema}[Tiga kasus]
Misalkan $\bigl(A\mid\mathbf{b}\bigr)$ berekuivalen baris dengan bentuk eselon tereduksi
$\bigl(R\mid\mathbf{c}\bigr)$, dengan $A$ berukuran $m\times n$.
\begin{enumerate}
  \item Sistem \emph{tidak konsisten} tepat ketika ada pivot pada kolom $(n+1)$.
  \item Bila konsisten, sistem punya \emph{jawaban tunggal} tepat ketika setiap kolom $A$
        memuat pivot.
  \item Bila konsisten, sistem punya \emph{banyak jawaban} tepat ketika ada kolom $A$
        tanpa pivot, dan banyaknya parameter bebas sama dengan banyaknya kolom tanpa pivot.
\end{enumerate}
\end{teorema}

\begin{proof}
Pivot pada kolom terakhir berarti ada baris berbentuk $0=\text{konstanta tak nol}$, yang
mustahil dipenuhi. Bila tidak ada pivot semacam itu, setiap baris tak nol memberi satu
persamaan yang menyatakan variabel pivot dalam variabel lain, sehingga sistem konsisten.

Bila setiap kolom memuat pivot, banyaknya pivot sama dengan $n$, sehingga tidak ada variabel
bebas dan setiap variabel pivot ditentukan tunggal. Bila ada kolom tanpa pivot, variabel
pada kolom itu bebas: ia boleh bernilai sebarang, dan setiap pilihan menentukan variabel
pivot secara tunggal. Jadi jawabannya tak berhingga banyak, dan banyaknya parameter bebas
tepat banyaknya kolom tanpa pivot.
\end{proof}

Gambar~\ref{fig:tiga-nasib} menaruh ketiga nasib itu pada satu bidang. Garis $x+y=2$
tidak berubah; yang berubah hanya garis kedua. Panel kiri: $x-y=0$ memotongnya di $(1,1)$,
dua pivot, jawaban tunggal. Panel tengah: $x+y=0$ sejajar, kolom terakhir matriks diperbesar
memuat pivot. Panel kanan: $2x+2y=4$ adalah garis yang sama dengan angka berbeda, satu pivot
dan satu parameter bebas.

\begin{figure}[htbp]
\centering
\begin{tikzpicture}[scale=0.66,>=Stealth]
  \begin{scope}
    \draw[layergrey!28,step=1] (-2.2,-2.2) grid (2.2,2.2);
    \draw[->,ink] (-2.3,0)--(2.4,0) node[right,font=\scriptsize]{$x$};
    \draw[->,ink] (0,-2.3)--(0,2.4) node[above,font=\scriptsize]{$y$};
    \draw[very thick,layerindigo] (-1.2,3.2)--(3.2,-1.2);
    \draw[very thick,layerteal] (-1.6,-1.6)--(2.2,2.2);
    \fill[accent] (1,1) circle (2.4pt);
    \node[font=\scriptsize,below right] at (1,1) {$(1,1)$};
    \node[font=\scriptsize] at (0,-2.9) {(a) berpotongan};
  \end{scope}
  \begin{scope}[xshift=6.4cm]
    \draw[layergrey!28,step=1] (-2.2,-2.2) grid (2.2,2.2);
    \draw[->,ink] (-2.3,0)--(2.4,0) node[right,font=\scriptsize]{$x$};
    \draw[->,ink] (0,-2.3)--(0,2.4) node[above,font=\scriptsize]{$y$};
    \draw[very thick,layerindigo] (-1.2,3.2)--(3.2,-1.2);
    \draw[very thick,accent,dashed] (-2.2,2.2)--(2.2,-2.2);
    \node[font=\scriptsize] at (0,-2.9) {(b) sejajar};
  \end{scope}
  \begin{scope}[xshift=12.8cm]
    \draw[layergrey!28,step=1] (-2.2,-2.2) grid (2.2,2.2);
    \draw[->,ink] (-2.3,0)--(2.4,0) node[right,font=\scriptsize]{$x$};
    \draw[->,ink] (0,-2.3)--(0,2.4) node[above,font=\scriptsize]{$y$};
    \draw[very thick,layerindigo] (-1.2,3.2)--(3.2,-1.2);
    \draw[line width=2.6pt,accent,densely dotted] (-1.2,3.2)--(3.2,-1.2);
    \node[font=\scriptsize] at (0,-2.9) {(c) berimpit};
  \end{scope}
\end{tikzpicture}
\caption{Tiga nasib sistem $2\times2$ dengan garis pertama $x+y=2$ tetap. Kiri: garis kedua
$x-y=0$ memotong di $(1,1)$, dua pivot, jawaban tunggal. Tengah: garis kedua $x+y=0$
sejajar, kolom ruas kanan memuat pivot, tidak ada jawaban. Kanan: garis kedua $2x+2y=4$
berimpit dengan garis pertama, satu pivot, tak berhingga jawaban.}
\label{fig:tiga-nasib}
\end{figure}

\begin{bacakeras}
Dua suara. Suara pertama menyebut keadaan, suara kedua menyebut akibatnya. Bacakan
bergantian, lalu bertukar peran.

Satu sistem. Dua garis. Tiga nasib.

Berpotongan: satu titik, satu jawaban.

Sejajar: tidak bertemu, tidak ada jawaban.

Berimpit: satu garis menyamar menjadi dua, jawaban tak berhingga.

Siapa yang memutuskan, sebelum kita menghitung apa pun? Pivot yang memutuskan.

Dua pivot, dua kolom: jawaban tunggal.

Pivot di kolom terakhir: ada baris yang berbunyi nol sama dengan bukan nol.

Ada kolom tanpa pivot: satu parameter bebas, jawaban berderet ke satu arah.

Ulangi baris terakhir bersama-sama: hitung pivot dulu, baru bicara jawaban.
\end{bacakeras}

Matriks dari Bab 7 memperlihatkan dua nasib sekaligus,
\[
  A=\begin{pmatrix}1&2&1\\ 2&5&3\\ 3&7&4\end{pmatrix},
  \qquad \text{baris ketiga}=\text{baris pertama}+\text{baris kedua}.
\]
Untuk $\mathbf{b}=(4,9,13)$ eliminasi memberi
\[
  \left(\begin{array}{ccc|c}1&2&1&4\\ 2&5&3&9\\ 3&7&4&13\end{array}\right)
  \longrightarrow
  \left(\begin{array}{ccc|c}1&0&-1&2\\ 0&1&1&1\\ 0&0&0&0\end{array}\right).
\]
Ada dua pivot dan tiga kolom, jadi ada satu variabel bebas. Jawabannya
$(2,1,0)+t\,(1,-1,1)$ untuk $t\in\mathbb{R}$, dan pemeriksaannya langsung:
$A(2,1,0)^{\top}=(4,9,13)^{\top}$ sedangkan $A(1,-1,1)^{\top}=\mathbf{0}$, jadi arah
pergeseran itu memang arah buta matriks dari Bab 7.

Untuk $\mathbf{b}=(4,9,14)$ baris terakhir menjadi $0=1$, dan sistem tidak konsisten.
Matriks koefisiennya sama persis; satu angka di ruas kanan memindahkan nasib sistem dari
``banyak jawaban'' ke ``tidak ada jawaban''. Karena itu pertanyaan tentang konsistensi tidak
dapat dijawab dari $A$ saja.

Kasus yang paling sering muncul dalam praktik adalah $m=n$. Di sana ketiga nasib itu
bertumpu pada satu bilangan, dan bilangan itulah yang kita cari di bagian berikut.

\begin{akibat}
Misalkan $A$ berukuran $n\times n$. Sistem $A\mathbf{x}=\mathbf{b}$ punya jawaban tunggal
untuk setiap $\mathbf{b}\in\mathbb{R}^{n}$ tepat ketika $\operatorname{rank}A=n$.
\end{akibat}

\begin{proof}
Bila $\operatorname{rank}A=n$, bentuk eselon tereduksi $A$ adalah $I_{n}$, sehingga
eliminasi pada $\bigl(A\mid\mathbf{b}\bigr)$ selalu berakhir dengan pivot pada setiap kolom
$A$ dan tanpa pivot pada kolom terakhir; teorema sebelumnya memberi jawaban tunggal.
Sebaliknya, bila $\operatorname{rank}A<n$, ada kolom tanpa pivot, dan untuk $\mathbf{b}$ yang
konsisten jawabannya tidak tunggal. Karena rank kurang dari $n$, ruang kolom $A$ tidak
memenuhi $\mathbb{R}^{n}$, sehingga ada $\mathbf{b}$ di luarnya; untuk $\mathbf{b}$ itu
sistemnya tidak konsisten.
\end{proof}

\section{Determinan sebagai fungsi bergantian}
\tingkat{menengah}

Kita ingin satu bilangan yang menjawab pertanyaan akibat sebelumnya tanpa menyelesaikan
sistem apa pun. Bilangan itu harus nol tepat ketika kolom-kolom $A$ tidak bebas linear,
berubah tanda ketika dua kolom bertukar, dan bernilai satu untuk matriks identitas. Ketiga
syarat itu ternyata sudah menentukan bilangannya secara tunggal.

\kamus{determinan}{determinant}{Faktor pengali volume: berapa kali ruang diregangkan atau
digepengkan oleh matriks itu.}

\begin{definisi}[Determinan]
Fungsi $D$ dari matriks $n\times n$ ke $\mathbb{R}$ disebut \emph{determinan} bila
\begin{enumerate}
  \item \emph{multilinear pada setiap baris}: untuk skalar $\alpha,\beta$ dan baris
        $\mathbf{u},\mathbf{v}$,
        $D(\dots,\alpha\mathbf{u}+\beta\mathbf{v},\dots)
         =\alpha D(\dots,\mathbf{u},\dots)+\beta D(\dots,\mathbf{v},\dots)$;
  \item \emph{bergantian}: $D(\dots,\mathbf{u},\dots,\mathbf{u},\dots)=0$ bila dua baris
        sama;
  \item \emph{ternormalisasi}: $D(I_{n})=1$.
\end{enumerate}
Nilai $D(A)$ ditulis $\det A$.\index{determinan}
\end{definisi}

Multilinearitas berarti determinan linear pada satu baris bila baris lain ditahan. Sifat
bergantian adalah syarat terkuat, dan akibatnya panjang.

\begin{lemma}
Menukar dua baris mengubah tanda determinan. Bila satu baris merupakan kelipatan skalar baris
lain, atau bila ada baris nol, determinannya nol.
\end{lemma}

\begin{proof}
Misalkan dua baris bersebelahan keduanya berisi $\mathbf{u}+\mathbf{v}$. Karena dua baris
sama, nilainya nol. Multilinearitas menguraikan nol itu menjadi empat suku,
\[
  0=D(\dots,\mathbf{u},\mathbf{u},\dots)+D(\dots,\mathbf{u},\mathbf{v},\dots)
   +D(\dots,\mathbf{v},\mathbf{u},\dots)+D(\dots,\mathbf{v},\mathbf{v},\dots),
\]
dan suku pertama serta keempat lenyap. Jadi menukar dua baris bersebelahan membalik tanda:
\[
  D(\dots,\mathbf{v},\mathbf{u},\dots)=-D(\dots,\mathbf{u},\mathbf{v},\dots).
\]
Penukaran yang tidak bersebelahan dicapai dengan beberapa penukaran bersebelahan, dan setiap
langkah membalik tanda.

Untuk kelipatan, $D(\dots,c\mathbf{u},\dots,\mathbf{u},\dots)
=c\,D(\dots,\mathbf{u},\dots,\mathbf{u},\dots)=0$ menurut multilinearitas dan sifat
bergantian. Baris nol adalah kasus $c=0$ dari bentuk yang sama.
\end{proof}

Sekarang ketunggalannya. Rumus yang muncul di tengah bukti adalah rumus Leibniz, dan rumus
itu nanti memberi ekspansi kofaktor sebagai akibat.

\begin{teorema}[Ketunggalan dan rumus Leibniz]
Hanya ada satu fungsi yang memenuhi ketiga syarat definisi, dan nilainya adalah
\[
  \det A=\sum_{\sigma}\operatorname{sgn}(\sigma)\,
  a_{\sigma(1)1}a_{\sigma(2)2}\cdots a_{\sigma(n)n},
\]
dengan jumlah diambil atas semua permutasi $\sigma$ dari $\{1,\dots,n\}$ dan
$\operatorname{sgn}(\sigma)=(-1)^{k}$ bila $\sigma$ dicapai dengan $k$ penukaran.
\end{teorema}

\begin{proof}
Tulis baris ke-$i$ sebagai $\mathbf{a}_{i}=\sum_{j}a_{ij}\mathbf{e}_{j}$. Multilinearitas
memungkinkan mengeluarkan setiap baris menjadi jumlah,
\[
  D(A)=\sum_{j_{1},\dots,j_{n}}a_{1j_{1}}a_{2j_{2}}\cdots a_{nj_{n}}\,
  D(\mathbf{e}_{j_{1}},\dots,\mathbf{e}_{j_{n}}).
\]
Suku dengan $j_{p}=j_{q}$ untuk $p\neq q$ lenyap karena $D$ bergantian, jadi hanya
$(j_{1},\dots,j_{n})$ yang merupakan permutasi $\sigma$ yang bertahan. Menurut lemma
sebelumnya, $D(\mathbf{e}_{\sigma(1)},\dots,\mathbf{e}_{\sigma(n)})
=\operatorname{sgn}(\sigma)D(I_{n})=\operatorname{sgn}(\sigma)$.

Untuk ketunggalan, misalkan $D$ memenuhi ketiga syarat. Perhitungan yang sama memberi
$D(A)=\bigl(\sum_{\sigma}\operatorname{sgn}(\sigma)a_{\sigma(1)1}\cdots
a_{\sigma(n)n}\bigr)D(I_{n})=\det A$. Jadi setiap fungsi yang memenuhi ketiga syarat sama
dengan rumus itu.
\end{proof}

Rumus Leibniz benar tetapi tidak praktis: ia menjumlahkan $n!$ suku, dan untuk $n=12$
banyaknya $479\,001\,600$. Yang dipakai orang adalah bentuk lain dari rumus yang sama.

\begin{akibat}[Ekspansi kofaktor]
Untuk setiap baris $i$ yang tetap berlaku
\[
  \det A=\sum_{j=1}^{n}(-1)^{i+j}a_{ij}M_{ij},
\]
dengan $M_{ij}$ adalah determinan submatriks $(n-1)\times(n-1)$ hasil menghapus baris $i$
dan kolom $j$. Angka $C_{ij}=(-1)^{i+j}M_{ij}$ disebut \emph{kofaktor} entri $a_{ij}$.
Rumus yang sama berlaku bila ekspansi dilakukan sepanjang kolom.
\end{akibat}

\begin{proof}
Dalam rumus Leibniz, kelompokkan suku-suku menurut nilai $j$ pada posisi baris $i$. Setiap
suku yang memuat $a_{ij}$ berbentuk $a_{ij}$ kali tanda kali hasilkali $n-1$ entri lain,
satu dari setiap baris selain $i$ dan satu dari setiap kolom selain $j$. Hasilkali itu tepat
suku Leibniz untuk minor $M_{ij}$, dan tandanya $(-1)^{i+j}$ karena memindahkan baris $i$
ke posisi terakhir menuntut $n-i$ penukaran dan memindahkan kolom $j$ ke posisi terakhir
menuntut $n-j$ penukaran, sehingga tandanya $(-1)^{2n-i-j}=(-1)^{i+j}$. Penjumlahan atas
$j$ memberi rumus yang ditulis.

Untuk ekspansi sepanjang kolom, dipakai $\det A=\det A^{\top}$. Kesamaan itu sendiri
mengikuti rumus Leibniz: setiap suku untuk $A^{\top}$ berbentuk
$a_{\sigma(1)1}\cdots a_{\sigma(n)n}$ dengan urutan indeks terbalik, dan pemetaan
$\sigma\mapsto\sigma^{-1}$ mempertahankan tanda karena banyaknya penukaran sama.
\end{proof}

\begin{contoh}
Untuk
$A=\begin{pmatrix}2&1&-1\\ -3&-1&2\\ -2&1&2\end{pmatrix}$ ekspansi sepanjang baris pertama
memberi kofaktor
\[
  C_{11}=\begin{vmatrix}-1&2\\ 1&2\end{vmatrix}=-4,\qquad
  C_{12}=-\begin{vmatrix}-3&2\\ -2&2\end{vmatrix}=2,\qquad
  C_{13}=\begin{vmatrix}-3&-1\\ -2&1\end{vmatrix}=-5,
\]
sehingga $\det A=2(-4)+1(2)+(-1)(-5)=-8+2+5=-1$. Tanda bergantian itu bukan hiasan: tanpa
tanda, hasilnya $-8+2-5=-11$, dan semua kesimpulan sesudahnya akan salah.
\end{contoh}

Determinan berguna untuk matriks hasil kali karena ia multiplikatif.

\begin{teorema}
Untuk matriks $n\times n$ berlaku $\det(AB)=\det A\,\det B$.
\end{teorema}

\begin{proof}
Tetapkan $A$ dan definisikan $f(B)=\det(AB)$. Baris ke-$i$ dari $AB$ adalah kombinasi linear
baris-baris $B$ dengan koefisien dari baris ke-$i$ $A$, sehingga $f$ multilinear terhadap
baris-baris $B$. Bila dua baris $B$ sama, dua baris $AB$ juga sama, jadi $f$ bergantian.
Menurut ketunggalan, setiap fungsi bergantian multilinear $f$ memenuhi
$f(B)=\det(B)\,f(I_{n})$. Di sini $f(I_{n})=\det A$, jadi
$f(B)=\det B\,\det A$.
\end{proof}

\begin{akibat}
$A$ invertibel jika dan hanya jika $\det A\neq0$. Bila invertibel,
$\det(A^{-1})=1/\det A$. Selain itu $\det(cA)=c^{n}\det A$, dan determinan matriks
segitiga adalah hasilkali entri diagonalnya.
\end{akibat}

\begin{proof}
Bila $A$ invertibel, $\det A\,\det(A^{-1})=\det I_{n}=1$, sehingga $\det A\neq0$ dan
inversnya $1/\det A$. Bila $A$ tidak invertibel, kolom-kolomnya bergantung linear menurut
Bab 7, sehingga salah satu kolom merupakan kombinasi linear kolom lain. Multilinearitas
menguraikan $\det A$ menjadi jumlah determinan dengan dua kolom sama, jadi semuanya nol.
Rumus $\det(cA)=c^{n}\det A$ diperoleh dengan mengeluarkan faktor $c$ dari setiap baris.
Untuk matriks segitiga, ekspansi kofaktor sepanjang baris yang hanya bersisi satu entri
diagonal, diulang $n$ kali, menyisakan hasilkali diagonal.
\end{proof}

Akibat itu membuat eliminasi dan determinan bertemu. Eliminasi membawa $A$ ke matriks
segitiga atas $U$ lewat penukaran dan penambahan kelipatan baris. Penambahan tidak mengubah
determinan; penukaran membalik tandanya. Karena itu $\det A=(-1)^{s}\det U$ dengan $s$
banyaknya penukaran, dan $\det U$ adalah hasilkali pivot: determinan gratis setelah eliminasi,
dan begitulah pustaka menghitungnya.

Sisi lain determinan yang sama pentingnya adalah luas. Kolom-kolom
$B=\begin{pmatrix}3&1\\ 1&2\end{pmatrix}$ adalah dua anak panah, $(3,1)$ dan $(1,2)$.
Persegi satuan yang sisinya dibawa oleh kedua panah itu berubah menjadi jajaran genjang,
dan luasnya lima, tepat $\det B=3\cdot2-1\cdot1=5$. Gambar~\ref{fig:luas-determinan}
memperlihatkan pemetaan itu. Kalau kedua kolom segaris, jajaran genjangnya gepeng menjadi
satu ruas, luasnya nol, dan determinannya juga nol.

\begin{figure}[htbp]
\centering
\begin{tikzpicture}[scale=0.95,>=Stealth]
  \draw[layergrey!28,step=1] (-0.4,-0.4) grid (4.4,3.4);
  \draw[->,ink] (-0.5,0)--(4.6,0) node[right,font=\scriptsize]{$x$};
  \draw[->,ink] (0,-0.5)--(0,3.5) node[above,font=\scriptsize]{$y$};
  \fill[layerteal!16] (0,0)--(3,1)--(4,3)--(1,2)--cycle;
  \draw[thick,layerteal] (0,0)--(3,1)--(4,3)--(1,2)--cycle;
  \fill[layerorange!22] (0,0)--(1,0)--(1,1)--(0,1)--cycle;
  \draw[thick,layerorange] (0,0)--(1,0)--(1,1)--(0,1)--cycle;
  \draw[very thick,accent,->] (0,0)--(3,1)
    node[below right,font=\scriptsize]{$(3,1)$};
  \draw[very thick,accent,->] (0,0)--(1,2)
    node[above left,font=\scriptsize]{$(1,2)$};
  \foreach \t in {1,2,3,4} {\node[font=\tiny,below] at (\t,-0.02) {\t};}
  \foreach \t in {1,2,3} {\node[font=\tiny,left] at (-0.02,\t) {\t};}
  \node[font=\tiny,fill=white,inner sep=0.5pt] at (0.5,0.5) {1};
  \node[font=\scriptsize,fill=white,inner sep=1.5pt] at (2,1.5) {luas $5$};
\end{tikzpicture}
\caption{Determinan sebagai luas. Persegi satuan (jingga, luas $1$) dipetakan oleh
$B=\begin{pmatrix}3&1\\ 1&2\end{pmatrix}$ menjadi jajaran genjang (teal) yang dibangun
$(3,1)$ dan $(1,2)$. Luasnya $5$, sama dengan $\det B=3\cdot2-1\cdot1=5$: determinan
mengalikan luas sebesar lima, dan tandanya mencatat apakah arah putaran dibalik.}
\label{fig:luas-determinan}
\end{figure}

\begin{tangan}
Bahan: kertas berpetak, gunting, penggaris, lem.
\begin{enumerate}
  \item Gambar persegi panjang tiga petak kali satu petak. Luasnya tiga petak.
  \item Gunting persegi panjang itu. Dari ujung kiri, potong satu segitiga siku-siku dengan alas
        dua petak dan tinggi satu petak, lalu tempelkan potongan itu ke ujung kanan.
  \item Bentuk barunya jajaran genjang yang miring. Luasnya tetap tiga petak, karena tidak
        ada kertas yang dibuang dan tidak ada yang ditambah.
  \item Ulangi dengan kemiringan satu petak, lalu dua, lalu tiga. Setiap kali hitung alas
        kali tinggi, jangan panjang sisi miring.
\end{enumerate}
Yang harus pembaca lihat: bentuk berubah, luas tidak. Kemiringan hanya menggeser, tidak
meregangkan, jadi determinan tidak bertambah.
\end{tangan}

\begin{untukanak}
Ambil selembar kertas persegi panjang, lalu miringkan dengan mendorong sisi atasnya ke samping
tanpa memendekkan atau memanjangkan apa pun. Tingginya tetap, lebarnya tetap, luasnya tetap,
walaupun bentuknya bukan persegi panjang lagi. Yang berubah hanya arah putaran: telusuri sisi
bawah lalu sisi kanan dari pojok kiri bawah, dan kertas miring masih berputar ke arah yang
sama; tukar kedua sisi, dan arahnya berbalik seperti sepatu kiri di kaki kanan.

Di bab ini kita akan tahu mengapa satu bilangan dapat memberitahu luas sekaligus arah, dan
mengapa bilangan itu menjadi nol tepat ketika kertas terlipat menjadi garis. Menurutmu, kalau
kertas itu dilipat dua, masih adakah cara menghitung luasnya dengan bilangan yang sama?
\end{untukanak}

\section{Aturan Cramer dan adjoin}
\tingkat{menengah}

Determinan menjawab apakah jawaban tunggal ada. Pertanyaan berikutnya: dapatkah jawaban itu
ditulis sebagai rumus, bukan sebagai hasil langkah? Untuk sistem kecil, dapat, dan rumusnya
rapi.

\kamus{adjoin}{adjugate}{Peta kofaktor yang disusun transpos, dipakai untuk menulis invers
sebagai satu rumus.}

\begin{definisi}[Adjoin]
Untuk $A$ berukuran $n\times n$ dengan kofaktor $C_{ij}$, \emph{matriks kofaktor} adalah
matriks dengan entri $C_{ij}$ pada posisi $(i,j)$, dan \emph{adjoin}
$\operatorname{adj}A$ adalah transposnya, yaitu matriks dengan entri $C_{ji}$ pada posisi
$(i,j)$.\index{kofaktor}\index{adjoin}
\end{definisi}

Perhatikan pembalikan indeksnya: kofaktor $C_{ij}$ berasal dari menghapus baris $i$ dan kolom
$j$, tetapi diletakkan pada posisi $(j,i)$. Inilah sumber kekeliruan paling sering saat
menghitung adjoin dengan tangan.

\begin{teorema}
Untuk setiap matriks $n\times n$ berlaku $A\operatorname{adj}A
=(\det A)I_{n}=\operatorname{adj}A\,A$.
\end{teorema}

\begin{proof}
Entri $(i,j)$ dari $A\operatorname{adj}A$ adalah
$\sum_{k}a_{ik}C_{jk}$. Bila $i=j$, jumlah itu tepat ekspansi kofaktor sepanjang baris $i$,
yaitu $\det A$. Bila $i\neq j$, jumlah itu adalah ekspansi sepanjang baris $j$ dari matriks
yang diperoleh dari $A$ dengan mengganti baris $j$ oleh baris $i$. Matriks itu punya dua
baris sama, jadi determinannya nol. Karena itu $A\operatorname{adj}A=(\det A)I_{n}$.
Kesamaan kedua diperoleh dengan argumen yang sama pada kolom, atau dari
$\operatorname{adj}(A^{\top})=(\operatorname{adj}A)^{\top}$.
\end{proof}

\begin{akibat}
Bila $\det A\neq0$, maka $A^{-1}=\dfrac{1}{\det A}\operatorname{adj}A$.
\end{akibat}

\begin{proof}
Dari teorema sebelumnya, $A\bigl(\operatorname{adj}A/\det A\bigr)=I_{n}$ dan
$\bigl(\operatorname{adj}A/\det A\bigr)A=I_{n}$. Invers matriks tunggal, jadi matriks itu
adalah $A^{-1}$.
\end{proof}

Rumus itu berlaku untuk semua $n$, tetapi harganya mahal: $\operatorname{adj}A$ memuat
$n^{2}$ kofaktor, dan setiap kofaktor adalah determinan $(n-1)\times(n-1)$.

Dengan rumus invers itu, jawaban sistem dapat ditulis sebagai nisbah dua determinan.

\begin{teorema}[Aturan Cramer]
Misalkan $A$ berukuran $n\times n$ dengan $\det A\neq0$, dan $\mathbf{b}\in\mathbb{R}^{n}$.
Jawaban tunggal $A\mathbf{x}=\mathbf{b}$ adalah
\[
  x_{j}=\frac{\det A_{j}}{\det A},\qquad j=1,\dots,n,
\]
dengan $A_{j}$ adalah matriks yang diperoleh dari $A$ dengan mengganti kolom ke-$j$ oleh
$\mathbf{b}$.\index{cramer}
\end{teorema}

\begin{proof}
Karena $\det A\neq0$, akibat sebelumnya memberi
$\mathbf{x}=A^{-1}\mathbf{b}=(\operatorname{adj}A)\,\mathbf{b}/\det A$. Entri ke-$j$ dari
hasil kali itu adalah $\sum_{i}b_{i}C_{ji}/\det A$. Jumlah $\sum_{i}b_{i}C_{ji}$ adalah
ekspansi kofaktor sepanjang kolom ke-$j$ dari matriks yang kolom ke-$j$-nya telah diganti
oleh $\mathbf{b}$, yaitu $A_{j}$. Jadi $x_{j}=\det A_{j}/\det A$.
\end{proof}

\begin{contoh}
Kembali ke sistem pembuka, $A\mathbf{x}=\mathbf{b}$ dengan
$\mathbf{b}=(8,-11,-3)^{\top}$. Kita sudah menghitung $\det A=-1$, dan
\[
  \det A_{1}=\begin{vmatrix}8&1&-1\\ -11&-1&2\\ -3&1&2\end{vmatrix}=-2,\qquad
  \det A_{2}=-3,\qquad \det A_{3}=1 .
\]
Jadi $x_{1}=(-2)/(-1)=2$, $x_{2}=(-3)/(-1)=3$, dan $x_{3}=1/(-1)=-1$. Pemeriksaan:
$A(2,3,-1)^{\top}=(4+3+1,\,-6-3-2,\,-4+3-2)^{\top}=(8,-11,-3)^{\top}$.

Adjoin matriks itu, dengan kofaktor baris pertamanya sudah kita hitung di atas, adalah
\[
  \operatorname{adj}A=\begin{pmatrix}-4&-3&1\\ 2&2&-1\\ -5&-4&1\end{pmatrix},
  \qquad
  A^{-1}=\frac{1}{-1}\operatorname{adj}A
  =\begin{pmatrix}4&3&-1\\ -2&-2&1\\ 5&4&-1\end{pmatrix}.
\]
Baris pertama $A^{-1}$ kali $\mathbf{b}$ memberi $4(8)+3(-11)-1(-3)=32-33+3=2$, cocok
dengan $x_{1}$.
\end{contoh}

Aturan Cramer indah dan jarang dipakai untuk menghitung. Nilainya terletak pada hal lain:
rumus itu menunjukkan jawaban bergantung secara mulus pada $\mathbf{b}$ dan pada entri $A$,
dan ketergantungan itulah yang kita ukur lewat biaya dan kepekaan.

\section{Matriks elementer, invers, dan syarat setara}
\tingkat{lanjut}

Eliminasi baris dapat dipandang sebagai perkalian matriks. Pandangan itu mengubah pertanyaan
``apakah $A$ invertibel'' menjadi pertanyaan tentang daftar faktor, dan daftar faktor itu
memberi algoritme invers.

\kamus{matriks elementer}{elementary matrix}{Matriks identitas yang satu langkahnya sudah
dikerjakan, sehingga menjadi catatan satu operasi baris.}

\begin{definisi}[Matriks elementer]
\emph{Matriks elementer} adalah matriks yang diperoleh dari $I_{n}$ dengan satu operasi
baris elementer. Ada tiga jenis, sesuai tiga operasi: penukaran dua baris, perkalian satu
baris dengan skalar tak nol, dan penambahan kelipatan satu baris pada baris lain.
\end{definisi}

\begin{lemma}
Bila $E$ matriks elementer $m\times m$ dan $A$ matriks $m\times n$, maka $EA$ adalah hasil
mengerjakan operasi baris yang sama pada $A$. Setiap matriks elementer invertibel, dan
inversnya matriks elementer sejenis.
\end{lemma}

\begin{proof}
Baris ke-$i$ dari $EA$ adalah kombinasi linear baris-baris $A$ dengan koefisien baris ke-$i$
$E$. Untuk $E$ jenis ketiga, baris itu adalah $\mathbf{a}_{i}+c\mathbf{a}_{j}$, tepat seperti
operasi barisnya; untuk dua jenis lain pemeriksaannya sama langsung. Inversnya diperoleh
dengan membatalkan operasi: menukar baris membatalkan dirinya sendiri, dan
$\operatorname{diag}(1,\dots,c,\dots,1)^{-1}=\operatorname{diag}(1,\dots,1/c,\dots,1)$.
\end{proof}

Gambar~\ref{fig:tiga-gerakan} menjalankan ketiga operasi pada sistem $x+2y=4$,
$3x+4y=10$, yang berpotongan di $(2,1)$. Menukar baris hanya menukar nama; mengalikan baris
pertama dengan $2$ menulis garis yang sama dengan angka dua kali; menambahkan baris kedua pada
baris pertama memberi $4x+6y=14$, yang masih melewati $(2,1)$. Titik potongnya tidak bergeser,
tetapi determinan matriks koefisiennya berganti tanda, berlipat dua, lalu tetap. Himpunan
jawaban kekal; bilangan yang mengukur luas tidak.

\begin{figure}[htbp]
\centering
\begin{tikzpicture}[scale=0.58,>=Stealth]
  \begin{scope}
    \draw[layergrey!28,step=1] (-0.4,-0.4) grid (4.4,3.2);
    \draw[->,ink] (-0.5,0)--(4.6,0) node[right,font=\scriptsize]{$x$};
    \draw[->,ink] (0,-0.5)--(0,3.4) node[above,font=\scriptsize]{$y$};
    \draw[very thick,layerindigo] (0,2)--(4,0) node[pos=0.15,above,font=\scriptsize]{$R_{1}$};
    \draw[very thick,layerteal] (0,2.5)--(3.33,0) node[pos=0.85,above left,font=\scriptsize]{$R_{2}$};
    \fill[accent] (2,1) circle (2.6pt);
    \node[font=\scriptsize,below right] at (2,1) {$(2,1)$};
    \node[font=\scriptsize] at (2,-1.1) {(i) tukar baris};
  \end{scope}
  \begin{scope}[xshift=5.6cm]
    \draw[layergrey!28,step=1] (-0.4,-0.4) grid (4.4,3.2);
    \draw[->,ink] (-0.5,0)--(4.6,0) node[right,font=\scriptsize]{$x$};
    \draw[->,ink] (0,-0.5)--(0,3.4) node[above,font=\scriptsize]{$y$};
    \draw[very thick,layerindigo] (0,2)--(4,0) node[pos=0.15,above,font=\scriptsize]{$R_{1}$};
    \draw[very thick,layerindigo,densely dashed] (0,2)--(4,0);
    \draw[very thick,layerteal] (0,2.5)--(3.33,0) node[pos=0.85,above left,font=\scriptsize]{$R_{2}$};
    \fill[accent] (2,1) circle (2.6pt);
    \node[font=\scriptsize,below right] at (2,1) {$(2,1)$};
    \node[font=\scriptsize] at (2,-1.1) {(ii) kalikan baris};
  \end{scope}
  \begin{scope}[xshift=11.2cm]
    \draw[layergrey!28,step=1] (-0.4,-0.4) grid (4.4,3.2);
    \draw[->,ink] (-0.5,0)--(4.6,0) node[right,font=\scriptsize]{$x$};
    \draw[->,ink] (0,-0.5)--(0,3.4) node[above,font=\scriptsize]{$y$};
    \draw[very thick,layergrey,dash pattern=on 3pt off 2pt] (0,2)--(4,0);
    \draw[very thick,layerteal] (0,2.5)--(3.33,0)
      node[pos=0.8,below right,font=\scriptsize]{$R_{2}$};
    \draw[very thick,layerplum] (0,2.33)--(3.5,0)
      node[pos=0.2,above,font=\scriptsize]{$R_{1}+R_{2}$};
    \fill[accent] (2,1) circle (2.6pt);
    \node[font=\scriptsize,below right] at (2,1) {$(2,1)$};
    \node[font=\scriptsize] at (2,-1.1) {(iii) tambah kelipatan};
  \end{scope}
\end{tikzpicture}
\caption{Tiga operasi baris elementer pada sistem $x+2y=4$, $3x+4y=10$. Menukar baris
hanya menukar nama $R_{1}$ dan $R_{2}$; mengalikan $R_{1}$ dengan $2$ menulis garis yang
sama; menambahkan $R_{2}$ pada $R_{1}$ memutar garis pertama menjadi $4x+6y=14$. Pada ketiga
panel, kedua garis melewati titik potong yang sama, $(2,1)$.}
\label{fig:tiga-gerakan}
\end{figure}

\begin{teorema}
Matriks $A$ berukuran $n\times n$ invertibel jika dan hanya jika $A$ merupakan hasilkali
matriks-matriks elementer.
\end{teorema}

\begin{proof}
Hasilkali matriks-matriks invertibel juga invertibel, jadi satu arah jelas.

Untuk arah lain, misalkan $A$ invertibel dan bawa $A$ ke bentuk eselon tereduksi lewat
operasi baris, yang sama dengan mengalikan dari kiri oleh matriks elementer
$E_{1},\dots,E_{k}$. Menurut teorema tiga kasus, $A$ invertibel berarti setiap kolom $A$
memuat pivot, sehingga bentuk tereduksi $A$ adalah $I_{n}$, sebab $A$ persegi dengan $n$
pivot. Jadi $E_{k}\cdots E_{1}A=I_{n}$, dan karena itu
$A=E_{1}^{-1}\cdots E_{k}^{-1}$, hasilkali matriks elementer.
\end{proof}

Bukti itu sekaligus algoritme: mengerjakan operasi yang sama pada $I_{n}$ memberi $A^{-1}$,
jadi $\bigl(A\mid I_{n}\bigr)\to\bigl(I_{n}\mid A^{-1}\bigr)$ menyelesaikan invers tanpa
satu pun determinan, jauh lebih murah daripada adjoin.

Kita kumpulkan yang sudah terbukti menjadi satu daftar. Daftar inilah yang dipakai orang
setiap hari, dan isinya wajah lain dari satu pernyataan.

\begin{teorema}[Syarat setara]
Untuk matriks $A$ berukuran $n\times n$, pernyataan berikut setara.
\begin{enumerate}
  \item $A$ invertibel.
  \item $A\mathbf{x}=\mathbf{0}$ hanya punya jawaban trivial $\mathbf{x}=\mathbf{0}$.
  \item Bentuk eselon tereduksi $A$ adalah $I_{n}$.
  \item $A$ adalah hasilkali matriks-matriks elementer.
  \item $A\mathbf{x}=\mathbf{b}$ konsisten untuk setiap $\mathbf{b}\in\mathbb{R}^{n}$.
  \item $A\mathbf{x}=\mathbf{b}$ punya tepat satu jawaban untuk setiap
        $\mathbf{b}\in\mathbb{R}^{n}$.
  \item $\det A\neq0$.
\end{enumerate}
\end{teorema}

\begin{proof}
Kita rangkaikan implikasi $(1)\Rightarrow(2)\Rightarrow(3)\Rightarrow(4)\Rightarrow(1)$
lalu $(1)\Rightarrow(5)\Rightarrow(6)\Rightarrow(1)$ dan $(1)\Leftrightarrow(7)$.

$(1)\Rightarrow(2)$: kalikan $A\mathbf{x}=\mathbf{0}$ dari kiri dengan $A^{-1}$.
$(2)\Rightarrow(3)$: inti $A$ hanya $\{\mathbf{0}\}$ berarti nullity nol, sehingga menurut
rank-nullity di Bab 7, $\operatorname{rank}A=n$; jadi ada $n$ pivot dan bentuk tereduksinya
$I_{n}$. $(3)\Rightarrow(4)$: langkah eliminasinya memberi
$E_{k}\cdots E_{1}A=I_{n}$. $(4)\Rightarrow(1)$: hasilkali matriks-matriks invertibel juga
invertibel.
$(1)\Rightarrow(5)$: ambil $\mathbf{x}=A^{-1}\mathbf{b}$. $(5)\Rightarrow(6)$: misalkan ada
$\mathbf{b}$ dengan dua jawaban berbeda $\mathbf{x}\neq\mathbf{y}$. Maka
$\mathbf{x}-\mathbf{y}\neq\mathbf{0}$ dan $A(\mathbf{x}-\mathbf{y})=\mathbf{0}$, sehingga
kolom-kolom $A$ bergantung linear dan $\operatorname{rank}A<n$. Ruang kolom $A$ karena itu
berdimensi kurang dari $n$, sehingga ada $\mathbf{b}'$ yang tidak terjangkau, bertentangan
dengan (5).
$(6)\Rightarrow(1)$: bila $A$ tidak invertibel, kolom-kolomnya bergantung, jadi ada
$\mathbf{x}\neq\mathbf{0}$ dengan $A\mathbf{x}=\mathbf{0}$, sehingga jawaban untuk
$\mathbf{b}=\mathbf{0}$ tidak tunggal. Terakhir $(1)\Leftrightarrow(7)$ adalah akibat di
bagian sebelumnya.
\end{proof}

Daftar itu menjelaskan mengapa empat pertanyaan yang tampak berbeda, tentang solusi, bentuk
eselon, faktorisasi, dan determinan, sebenarnya satu pertanyaan; yang berubah hanya alat
ukurnya. Tetapi daftar itu diam tentang biaya dan kepekaan, dan dua hal itulah yang menentukan
apakah jawaban yang diperoleh dapat dipercaya.

\begin{tantangan}
Sistem berikut tampak biasa: tiga persamaan untuk tiga bilangan, tidak ada baris kembar yang
mencolok, sehingga kelihatannya jawabannya tunggal.
\[
  x+y+z=3,\qquad x+2y+3z=6,\qquad 2x+3y+4z=9 .
\]
Tentukan semua jawabannya. Berapa banyak jawabannya, dan mengapa menghitung banyaknya
persamaan menyesatkan di sini?

Petunjuk 1: jumlahkan dua persamaan pertama, lalu bandingkan hasilnya dengan persamaan
ketiga.

Petunjuk 2: persamaan ketiga tidak menambah keterangan apa pun, jadi hanya dua keterangan
yang bebas. Setelah $z$ dipilih, $y$ lalu $x$ ikut tertentu; tuliskan jawaban sebagai satu
keluarga dengan satu parameter.

Kalau mentok, kembali ke bagian tiga nasib dan hitung pivotnya.
\end{tantangan}

\subsection*{Biaya dan kepekaan}

Eliminasi pada matriks $n\times n$ memakai sekitar $n^{3}/3$ perkalian: untuk $n=1000$
sekitar $3{,}33\times10^{8}$, yang masih cepat. Aturan Cramer dengan kofaktor memakai
$n\cdot n!$ perkalian, sekitar $5{,}75\times10^{9}$ untuk $n=12$ dan di atas $10^{19}$ untuk
$n=20$; karena itu ia hanya dipakai untuk $n$ kecil atau untuk keperluan teoretis.

Yang lebih halus adalah kepekaan. Misalkan $A$ invertibel dan $\mathbf{b}\neq\mathbf{0}$,
dan misalkan $\mathbf{b}$ berubah menjadi $\mathbf{b}+\delta\mathbf{b}$ sehingga jawabannya
berubah menjadi $\mathbf{x}+\delta\mathbf{x}$.

\begin{teorema}
Dengan notasi di atas, untuk sembarang norma matriks yang subordinate berlaku
\[
  \frac{\lVert\delta\mathbf{x}\rVert}{\lVert\mathbf{x}\rVert}
  \le\kappa(A)\,
  \frac{\lVert\delta\mathbf{b}\rVert}{\lVert\mathbf{b}\rVert},
  \qquad \kappa(A)=\lVert A\rVert\,\lVert A^{-1}\rVert .
\]
\end{teorema}

\begin{proof}
Dari $A\,\delta\mathbf{x}=\delta\mathbf{b}$ diperoleh
$\delta\mathbf{x}=A^{-1}\delta\mathbf{b}$, sehingga
$\lVert\delta\mathbf{x}\rVert\le\lVert A^{-1}\rVert\lVert\delta\mathbf{b}\rVert$.
Dari $\mathbf{b}=A\mathbf{x}$ diperoleh
$\lVert\mathbf{b}\rVert\le\lVert A\rVert\lVert\mathbf{x}\rVert$, yaitu
$1/\lVert\mathbf{x}\rVert\le\lVert A\rVert/\lVert\mathbf{b}\rVert$. Mengalikan kedua
ketaksamaan memberi hasil yang diminta.
\end{proof}

Bilangan $\kappa(A)$ disebut \emph{bilangan kondisi}. Ia tidak bergantung pada skala:
mengalikan $A$ dengan $c$ tidak mengubahnya, karena $\lVert cA\rVert\lVert (cA)^{-1}\rVert
=|c|\lVert A\rVert\cdot|c|^{-1}\lVert A^{-1}\rVert$. Determinan tidak punya sifat itu, dan
inilah alasan kita tidak memakai determinan sebagai ukuran kepekaan. Determinan kecil hanya
berarti volume kecil; ia dapat kecil karena seluruh matriks diperkecil, sedangkan kepekaan
tidak berubah.

Untuk sistem pembuka, $\kappa(S)\approx1{,}01\times10^{5}$. Dengan presisi mesin sekitar
$10^{-16}$, sisa angka yang dapat dipercaya pada jawaban sekitar $10^{-11}$ relatif: cukup
untuk contoh ini, habis untuk matriks yang lebih buruk. Bab 22 membahas ambang dan cara
mengenali matriks yang sudah tidak dapat dipercaya.
\index{kondisi}

\begin{untukanak}
Bayangkan papan panjang yang dipakai sebagai landasan untuk mengukur tinggi sebuah bukit.
Papan itu hampir mendatar, hanya miring sedikit. Kalau ujung yang di bawah diangkat sejari,
ujung yang lain naik jauh lebih banyak, karena papan itu bertumpu pada kemiringan yang
sangat kecil. Sejari di satu tempat menjadi sejengkal di tempat lain.

Sekarang bayangkan papannya dibuat sangat curam. Angkat ujung bawahnya sejari, dan ujung
atasnya hanya naik sedikit sekali. Pekerjaan yang sama menghasilkan dua akibat yang sangat
berbeda, hanya karena landasannya berbeda.

Hitungan kita bekerja seperti papan itu. Ada susunan angka yang membuat perubahan kecil pada
masukan menghasilkan perubahan kecil pada jawaban; ada pula susunan yang membuat perubahan
sebesar debu menghasilkan jawaban yang sama sekali lain. Yang menentukan bukan besar kecilnya
angka, melainkan bentuk susunannya.

Di bab ini kita akan tahu mengapa ada satu bilangan yang dapat meramalkan hal itu sebelum
kita menghitung, dan mengapa bilangan itu tidak berubah ketika semua angka diperbesar atau
diperkecil dengan perbandingan yang sama.

Kalau kamu menimbang dengan timbangan yang landasannya miring, apakah hasilnya salah, atau
hanya tidak dapat dipercaya?
\end{untukanak}

\section{Contoh hitung}
\tingkat{menengah}

Kita kerjakan satu sistem dengan dua cara, hitung satu determinan $4\times4$, lalu mengukur
kepekaan dua matriks.

\begin{resep}
\begin{enumerate}
  \item Tulis matriks diperbesar $\bigl(A\mid\mathbf{b}\bigr)$: kolom koefisien dulu, kolom
        ruas kanan paling akhir.
  \item Kerjakan kolom paling kiri. Pilih satu entri tak nol sebagai pivot; kalau entri
        teratas nol, tukar baris lebih dulu.
  \item Nolkan semua entri di bawah pivot: untuk tiap baris $j$ di bawahnya,
        $R_{j}\leftarrow R_{j}-\frac{a_{ji}}{a_{ii}}R_{i}$.
  \item Geser satu kolom ke kanan dan ulangi; berhenti saat tidak ada baris tak nol lagi,
        lalu baca nasib sistem dari letak pivot.
\end{enumerate}
Peringatan urutan: pivot dinolkan dari atas ke bawah, kolom demi kolom, sebelum substitusi
mundur dimulai. Membalik urutan itu menghapus pivot yang sudah diperoleh dan membuat
pekerjaan harus diulang.
\end{resep}

Sistemnya $A\mathbf{x}=\mathbf{b}$ dengan
\[
  A=\begin{pmatrix}2&1&-1\\ -3&-1&2\\ -2&1&2\end{pmatrix},
  \qquad \mathbf{b}=\begin{pmatrix}8\\ -11\\ -3\end{pmatrix}.
\]

Cara pertama, eliminasi, dengan langkah ditulis lengkap:
\[
  \left(\begin{array}{ccc|c}2&1&-1&8\\ -3&-1&2&-11\\ -2&1&2&-3\end{array}\right)
  \xrightarrow[R_{3}\leftarrow R_{3}+R_{1}]{R_{2}\leftarrow R_{2}+\tfrac{3}{2}R_{1}}
  \left(\begin{array}{ccc|c}2&1&-1&8\\ 0&\tfrac{1}{2}&\tfrac{1}{2}&1\\
  0&2&1&5\end{array}\right),
\]
lalu $R_{3}\leftarrow R_{3}-4R_{2}$ memberi
\[
  \left(\begin{array}{ccc|c}2&1&-1&8\\ 0&\tfrac{1}{2}&\tfrac{1}{2}&1\\
  0&0&-1&1\end{array}\right).
\]
Substitusi mundur: dari baris ketiga $-x_{3}=1$, jadi $x_{3}=-1$. Baris kedua
$\tfrac{1}{2}x_{2}+\tfrac{1}{2}(-1)=1$ memberi $x_{2}=3$. Baris pertama
$2x_{1}+3+1=8$ memberi $x_{1}=2$. Hasilnya $\mathbf{x}=(2,3,-1)$. Pivot yang muncul adalah
$2$, $\tfrac{1}{2}$, dan $-1$, jadi hasilkali pivot $\det A=2\cdot\tfrac{1}{2}\cdot(-1)=-1$,
cocok dengan ekspansi kofaktor tadi. Di sini tidak ada penukaran baris, sehingga tandanya
tidak berubah.

Cara kedua, aturan Cramer, sudah kita jalankan pada sistem yang sama di bagian sebelumnya.
Ketiga determinan pengganti memberi $\det A_{1}=-2$, $\det A_{2}=-3$, dan $\det A_{3}=1$,
sehingga $\mathbf{x}=(-2,-3,1)/(-1)=(2,3,-1)$. Dua jalan, satu jawaban; pemeriksaan semacam
ini murah dan sebaiknya selalu dikerjakan.

Sekarang determinan $4\times4$. Ambil
\[
  M=\begin{pmatrix}2&0&1&3\\ 1&4&0&2\\ 3&1&2&1\\ 0&2&1&5\end{pmatrix}.
\]
Ekspansi sepanjang baris pertama memakai empat minor, dan salah satunya kita hitung penuh:
\[
  M_{11}=\begin{vmatrix}4&0&2\\ 1&2&1\\ 2&1&5\end{vmatrix}
  =4(10-1)-0(5-2)+2(1-4)=36-6=30 .
\]
Dengan cara serupa $M_{12}=15$, $M_{13}=-45$, dan $M_{14}=-15$. Karena itu
\[
  \det M=2(30)-0(15)+1(-45)-3(-15)=60-45+45=60 .
\]
Banyaknya suku yang dijumlahkan adalah $4!=24$, dan untuk $5\times5$ sudah $120$:
pertumbuhannya faktorial. Inilah alasan pustaka tidak memakai cara ini.

Sekarang kepekaan. Matriks Hilbert $4\times4$ didefinisikan oleh $h_{ij}=1/(i+j-1)$.
Determinannya tepat $1/6\,048\,000\approx1{,}65\times10^{-7}$, sangat kecil, sehingga orang
mudah menyimpulkan bahwa matriks ini berbahaya. Bilangan kondisinya $15\,513{,}74$. Matriks
pembuka tadi punya $\det S=-0{,}003$ dan $\kappa(S)\approx1{,}01\times10^{5}$, jadi
determinannya delapan belas ribu kali lebih besar, sedangkan kepekaannya sekitar enam
setengah kali lebih buruk. Determinan bukan ukuran kepekaan; bilangan kondisi itu ukurannya.

Uji lompatan untuk $S$ memakai $\mathbf{x}=(1,1,1)$, sehingga
$\mathbf{b}=(6,\,15,\,24{,}001)$. Mengubah komponen ketiga menjadi $24{,}002$ mengubah
jawaban menjadi $(2,-1,2)$. Perubahan relatif data $3{,}46\times10^{-5}$, perubahan relatif
jawaban $\sqrt{2}\approx1{,}414$, nisbahnya $40\,916$: di bawah $\kappa(S)$ seperti
dijanjikan teorema, tetapi hanya dengan faktor dua setengah.

\kertas{Manakah yang lebih peka, $M$ atau matriks Hilbert $4\times4$? Ramalkan dengan membandingkan
determinannya, lalu hitung bilangan kondisi keduanya dengan \texttt{numpy.linalg.cond}.
Bandingkan juga dengan $10^{-3}I_{4}$, yang determinannya jauh lebih kecil, lalu jelaskan
mengapa aturan ``determinan kecil berarti peka'' gagal. Terakhir, kalikan $M$ dengan
$10^{-4}$: bilangan kondisinya tidak berubah, sedangkan determinannya mengecil sebesar
$10^{-16}$.}

\section{Lab komputasi}
\tingkat{lanjut}

\begin{lab}
\begin{verbatim}
import numpy as np, math

def det_kofaktor(A):
    n = A.shape[0]
    if n == 1:
        return A[0, 0]
    return sum((-1) ** j * A[0, j] *
               det_kofaktor(np.delete(np.delete(A, 0, 0), j, 1))
               for j in range(n))

A = np.array([[2., 1., -1.], [-3., -1., 2.], [-2., 1., 2.]])
b = np.array([8., -11., -3.])
x = np.linalg.solve(A, b)
print("det A    :", np.linalg.det(A))
print("det kof  :", det_kofaktor(A))
print("rank A   :", np.linalg.matrix_rank(A))
print("x        :", x)
print("sisa     :", np.linalg.norm(A @ x - b))
for j in range(3):
    Aj = A.copy(); Aj[:, j] = b
    print("det A%d   : %.1f  -> x%d = %.6f" % (j + 1, np.linalg.det(Aj),
          j + 1, np.linalg.det(Aj) / np.linalg.det(A)))

M = np.array([[2., 0., 1., 3.], [1., 4., 0., 2.],
              [3., 1., 2., 1.], [0., 2., 1., 5.]])
print("det M    :", det_kofaktor(M), np.linalg.det(M))

H = np.array([[1. / (i + j + 1) for j in range(4)] for i in range(4)])
print("det H4   :", np.linalg.det(H))
print("cond H4  : %.2f" % np.linalg.cond(H))

S = np.array([[1., 2., 3.], [4., 5., 6.], [7., 8., 9.001]])
bs = S @ np.array([1., 1., 1.])
print("det S    :", np.linalg.det(S), " cond S: %.1f" % np.linalg.cond(S))
print("x        :", np.linalg.solve(S, bs))
bs2 = bs.copy(); bs2[2] += 0.001
print("x'       :", np.linalg.solve(S, bs2))
print("biaya    :", math.factorial(12) * 12, round(1000 ** 3 / 3))
\end{verbatim}
Keluaran sebenarnya:
\begin{verbatim}
det A    : -1.0
det kof  : -1.0
rank A   : 3
x        : [ 2.  3. -1.]
sisa     : 1.9860273225978185e-15
det A1   : -2.0  -> x1 = 2.000000
det A2   : -3.0  -> x2 = 3.000000
det A3   : 1.0  -> x3 = -1.000000
det M    : 60.0 60.000000000000036
det H4   : 1.6534391534390412e-07
cond H4  : 15513.74
det S    : -0.0030000000000036655  cond S: 101072.3
x        : [1. 1. 1.]
x'       : [ 2. -1.  2.]
biaya    : 5748019200 333333333
\end{verbatim}
\end{lab}

Empat hal perlu dibaca dari keluaran itu.

Pertama, $x=(2,3,-1)$ muncul dengan sisa $1{,}99\times10^{-15}$, yaitu nol sampai presisi
mesin. Sisa itu sengaja dihitung: jawaban pustaka tidak pernah persis, dan sisa adalah cara
termurah memeriksa apakah jawaban itu memenuhi sistemnya. Pada matriks buruk, sisa kecil
tidak menjamin galat kecil; itu topik Bab 22.

Kedua, determinan hasil ekspansi kofaktor dan hasil pustaka sama-sama $-1$, dan untuk $M$
keduanya $60$. Pustaka tidak memakai kofaktor; ia memakai faktorisasi dan mengalikan pivot.
Selisih $3{,}6\times10^{-14}$ pada $\det M$ adalah pembulatan, bukan perbedaan definisi.

Ketiga, $\det H_{4}$ dilaporkan $1{,}65\times10^{-7}$ sedangkan nilainya tepat
$1/6\,048\,000$, jadi angka mesin itu benar sampai tujuh digit. Untuk orde lebih tinggi,
$1/\det$ melampaui $10^{12}$ dengan cepat, dan di situlah determinan berhenti berguna
sebagai penanda.

Keempat, lompatan jawaban pada $S$ terlihat utuh: dari $(1,1,1)$ ke $(2,-1,2)$ hanya karena
satu angka berubah seperseribu. Angka yang dilaporkan untuk $\det S$ pun tidak persis
$-0{,}003$; selisih $3{,}7\times10^{-15}$ muncul karena $9{,}001$ tidak terwakili tepat
dalam pecahan biner.

Terakhir, biaya: $5{,}748\times10^{9}$ perkalian untuk aturan Cramer pada $n=12$ melawan
$3{,}33\times10^{8}$ untuk eliminasi pada $n=1000$. Matriks yang jauh lebih besar dikerjakan
dengan pekerjaan yang jauh lebih sedikit; itu bedanya pangkat tiga dan faktorial.

\begin{untukdosen}
Urutan mengajar yang berhasil dipakai: eliminasi lebih dulu, dengan tiga kasus ditunjukkan
pada satu matriks koefisien tetap dan tiga ruas kanan berbeda. Mahasiswa yang melihat
konsistensi berubah karena satu angka di $\mathbf{b}$ biasanya langsung menerima bahwa
pertanyaan tentang $A$ dan pertanyaan tentang $\mathbf{b}$ harus dipisahkan. Baru setelah itu
determinan masuk sebagai jawaban atas pertanyaan tentang $A$ saja.

Dua miskonsepsi khas. Pertama, ``determinan kecil berarti sistem tidak stabil''. Matriks
$10^{-6}I$ punya determinan $10^{-18}$ dan bilangan kondisi satu; sebaliknya $H_{4}$ punya
determinan $1{,}65\times10^{-7}$ dan kondisi $1{,}5\times10^{4}$. Kekeliruan ini muncul
setiap tahun, dan cara membongkarnya adalah meminta mahasiswa mengalikan satu matriks
dengan skalar lalu membandingkan kedua bilangan. Kedua, ``aturan Cramer adalah cara
menyelesaikan sistem''. Secara teori benar; secara biaya salah. Beri satu sistem $12\times12$
dan minta mereka memperkirakan biaya kedua cara sebelum menghitung.

Aktivitas kelas yang berguna: bagikan tiga sistem dengan matriks koefisien sama dan ruas
kanan $(4,9,13)$, $(4,9,14)$, dan $(0,0,0)$. Minta setiap kelompok menentukan nasib
sistemnya hanya dari bentuk eselon, lalu menyimpulkan apa yang dapat dan tidak dapat
disimpulkan dari $A$ seorang diri.

Tanda kesiapan lanjut: mahasiswa dapat menyebutkan tujuh syarat setara tanpa menghafal
urutan, dan dapat menjelaskan mengapa syarat tentang determinan berada di daftar yang sama
dengan syarat tentang pivot.
\end{untukdosen}

\section{Ringkasan}
\tingkat{dasar}

\begin{ringkasan}
\begin{itemize}
  \item Operasi baris elementer tidak mengubah himpunan jawaban, karena setiap operasi dapat
        dibatalkan. Karena bentuk tereduksi tunggal, letak pivot menentukan tiga nasib: tidak
        konsisten bila pivot muncul di kolom terakhir, tunggal bila setiap kolom koefisien
        punya pivot, banyak bila ada kolom tanpa pivot.
  \item Determinan adalah satu-satunya fungsi multilinear, bergantian, dan ternormalisasi.
        Rumus Leibniz dan ekspansi kofaktor adalah akibatnya, bukan definisinya.
  \item $\det(AB)=\det A\det B$, sehingga $A$ invertibel tepat ketika $\det A\neq0$; dan
        $\det A$ sama dengan hasilkali pivot setelah eliminasi, kali tanda penukaran.
  \item Adjoin memberi $A\operatorname{adj}A=(\det A)I$ dan
        $A^{-1}=\operatorname{adj}A/\det A$; aturan Cramer menulis
        $x_{j}=\det A_{j}/\det A$.
  \item Tujuh syarat invertibilitas setara: inti trivial, bentuk tereduksi $I_{n}$, hasilkali
        matriks elementer, konsisten dan tunggal untuk semua $\mathbf{b}$, dan
        $\det A\neq0$.
  \item Biaya eliminasi sebanding $n^{3}/3$, biaya kofaktor sebanding $n\cdot n!$. Kepekaan
        diukur oleh $\kappa(A)=\lVert A\rVert\lVert A^{-1}\rVert$, bukan oleh $\det A$.
\end{itemize}
\end{ringkasan}

\section{Latihan}
\tingkat{menengah}

\begin{latihan}
\begin{enumerate}
  \item Selesaikan dengan eliminasi, lalu periksa dengan aturan Cramer:
        \[
          \begin{pmatrix}1&2&3\\ 2&5&3\\ 1&0&8\end{pmatrix}\mathbf{x}
          =\begin{pmatrix}9\\ 8\\ 26\end{pmatrix}.
        \]
        Tuliskan pivot yang muncul, catat berapa kali baris ditukar, lalu simpulkan
        $\det A$ dari hasilkali pivot.

  \item Hitung $\det M$ untuk
        $M=\begin{pmatrix}0&1&2&3\\ 1&0&3&2\\ 2&3&0&1\\ 3&2&1&0\end{pmatrix}$
        dengan ekspansi kofaktor sepanjang satu baris, lalu hitung dengan
        \texttt{numpy.linalg.det}. Kedua hasil tidak sama persis: yang pertama eksak, yang
        kedua $5{,}3\times10^{-15}$. Jelaskan mengapa. Sesudah itu cari vektor tak nol
        $\mathbf{v}$ dengan $M\mathbf{v}=\mathbf{0}$ dan simpulkan $\det M=0$ tanpa
        ekspansi apa pun.

  \item Buktikan bahwa bila dua kolom $A$ sama, maka $\det A=0$, tanpa memakai
        $\det A=\det A^{\top}$.
        (Petunjuk: pakai rumus Leibniz, kelompokkan suku-suku menurut pasangan permutasi yang
        bertukar pada dua kolom itu, dan tunjukkan tandanya berlawanan.)

  \item Buktikan bahwa $\det(\operatorname{adj}A)=(\det A)^{n-1}$ untuk $n\ge2$.
        (Petunjuk: ambil determinan pada $A\operatorname{adj}A=(\det A)I$ dan gunakan
        $\det(cI)=c^{n}$.)

  \item Nilai klaim ini: \emph{bila $\det A$ sangat kecil, maka jawaban $A\mathbf{x}
        =\mathbf{b}$ sangat peka terhadap perubahan $\mathbf{b}$}. Tentukan benar, salah,
        atau belum cukup bukti, dan sertakan contoh penyangkal atau alasan tambahan yang
        membuatnya benar.

  \item Untuk $A=\begin{pmatrix}1&1\\ 1&1+\varepsilon\end{pmatrix}$ dengan
        $\varepsilon>0$, hitung $\kappa(A)$ dalam norma $1$ sebagai fungsi
        $\varepsilon$. Tentukan nilai $\varepsilon$ yang membuat $\kappa(A)>10^{6}$, lalu
        periksa dengan \texttt{numpy.linalg.cond}.

  \item Buktikan: bila $A$ invertibel dan $B$ diperoleh dari $A$ dengan menambahkan
        kelipatan satu baris pada baris lain, maka $\det B=\det A$ dan
        $B^{-1}$ dapat dinyatakan lewat $A^{-1}$ dan satu matriks elementer.

  \item Misalkan $D=\operatorname{diag}(d_{1},\dots,d_{n})$ dan $A$ berukuran $n\times n$.
        Tunjukkan bahwa $\det(A+AD)=\det(A)\prod_{i=1}^{n}(1+d_{i})$, lalu tentukan syarat
        pada $d_{1},\dots,d_{n}$ agar $A+AD$ invertibel. Periksa jawabanmu dengan satu
        matriks $3\times3$ di Python.

  \item Pilih satu matriks $3\times3$ dengan entri bulat kecil dan $\det A=1$, hitung
        $A^{-1}$ lewat adjoin dan lewat \texttt{numpy.linalg.inv}, lalu bandingkan
        keduanya. Ulangi untuk matriks dengan $\det A=7$ dan jelaskan mengapa entri
        $A^{-1}$ menjadi pecahan.
\end{enumerate}

Petunjuk tiga butir tersulit. Butir 4: ambil determinan kedua ruas, ingat
$\det(cI)=c^{n}$, dan perhatikan bahwa $\det A$ boleh nol sehingga kesimpulan tentang
$\operatorname{adj}A$ harus dibaca sebagai matriks, bukan invers. Butir 6: dalam norma $1$,
$\lVert A\rVert_{1}$ adalah jumlah kolom terbesar, dan untuk matriks ini kolom keduanya
yang terbesar; hitung $\lVert A^{-1}\rVert_{1}$ dari rumus $2\times2$ dan perhatikan bahwa
jawabannya sebanding $1/\varepsilon$. Butir 8: tulis $A+AD$ sebagai $A(I+D)$ sebelum
mengeluarkan determinan, dan ingat bahwa $I+D$ bersifat diagonal sehingga determinannya
hasilkali entri diagonalnya.
\end{latihan}

\begin{untukpeneliti}
Yang benar-benar terbukti di bab ini, untuk medan $\mathbb{R}$ dan dimensi hingga: ketunggalan
determinan sebagai fungsi $n$-linear bergantian yang ternormalisasi beserta rumus permutasi,
di Hoffman dan Kunze, \emph{Linear Algebra} edisi ke-2, bagian 5.2 dan 5.3, berlaku lebih umum
atas gelanggang komutatif dengan identitas (bagian 5.1); multiplikativitas dan kesetaraan
invertibilitas dengan $\det\neq0$, di Axler, \emph{Linear Algebra Done Right} edisi ke-4,
bagian 9C, pernyataan 9.49 dan 9.50, dengan determinan operator yang bebas basis; serta daftar
tujuh syarat setara, di Anton edisi ke-11, Teorema 2.3.8.

Yang dipopulerkan: ``determinan adalah volume''. Yang terbukti adalah $|\det A|$ sama dengan
volume paralelepipedum kolom-kolomnya; tanpa nilai mutlak, klaim itu salah untuk determinan
negatif. Juga keliru secara numerik: aturan Cramer sebagai cara menyelesaikan sistem di atas
$n=11$, tempat biayanya faktorial.

Pertanyaan terbuka. Batas kepekaan pada teorema di atas bersifat normwise dan worst-case,
sering terlalu pesimistis dibanding analisis componentwise. Lebih mendasar: faktor pertumbuhan
maksimum eliminasi Gauss dengan pivot parsial belum diketahui untuk $n$ umum. Batas terbukti
$2^{n-1}$, sedangkan percobaan numerik memberi jauh lebih kecil, dan kesenjangan itu terbuka.

Batas keberlakuan. Argumen runtuh bila medan bukan lapangan, sebab modul tidak selalu punya
basis dan eliminasi tidak selalu dapat membagi. Untuk $n$ sangat besar, determinan meluap
atau mengalir habis jauh sebelum rank-nya hilang, sehingga pustaka melaporkan logaritma
determinan. Pernyataan kepekaan juga bergantung pada norma, dan tidak ada pilihan yang
kanonik.
\end{untukpeneliti}

\peta{Bab ini mengikuti Anton, \emph{Elementary Linear Algebra} edisi ke-11, bagian 1.1
sampai 1.6 (sistem linear, eliminasi Gauss-Jordan, invers, matriks elementer) dan 2.1 sampai
2.3 (kofaktor, determinan lewat reduksi baris, aturan Cramer, termasuk Teorema 2.3.8 tentang
syarat setara). Perlakuan aksiomatik mengikuti Hoffman dan Kunze, \emph{Linear Algebra} edisi
ke-2, bagian 5.2 dan 5.3, dengan bagian 1.4 untuk ketunggalan bentuk eselon tereduksi.
Pandangan bebas basis ada di Axler, \emph{Linear Algebra Done Right} edisi ke-4, bagian 9C.
Untuk eliminasi sebagai faktorisasi dan pandangan $A=CR$, lihat Strang, \emph{Introduction to
Linear Algebra} edisi ke-6, Bab 2 dan catatan ``Elimination and Factorization''. Sisi numerik,
yaitu bilangan kondisi, galat pembulatan, dan biaya, dibahas di Bab 22 buku ini.}
